\appendix
\section*{Anexos y apéndices}
\addcontentsline{toc}{section}{Anexos y apéndices}
\renewcommand{\thesubsection}{\Alph{subsection}}
En el \textbf{Anexo} se describe el material entregado por el docente (los datos de la instancia). Los \textbf{Apéndices} contienen material elaborado por los autores: resultados completos, la solución exacta, el código fuente y las instrucciones para reproducir todo.

\subsection*{Anexo A. Instancia entregada por el docente}
\addcontentsline{toc}{subsection}{Anexo A. Instancia entregada por el docente}

La instancia consiste en una matriz binaria $A$ de 500$\times$500 (filas: clientes; columnas: antenas), donde $a_{ij}=1$ si la antena $j$ cubre al cliente $i$, y un vector de costos $c_j$ por antena. La Tabla~\ref{tab:instancia} resume sus características.
\begin{table}[H]
\centering
\caption{Características de la instancia}
\label{tab:instancia}
\begin{tabular}{lr}
\toprule
Característica & Valor \\
\midrule
Clientes (restricciones) & 500 \\
Antenas (variables) & 500 \\
Densidad de la matriz (proporción de unos) & 0.100 \\
Antenas que cubren a cada cliente (mín. / media / máx.) & 32 / 49.9 / 70 \\
Clientes que cubre cada antena (mín. / media / máx.) & 33 / 49.9 / 66 \\
Costo por antena (mín. / media / máx.) & 2\,000 / 3\,034 / 3\,998 \\
Costo de instalar todas las antenas & 1\,516\,821 \\
\bottomrule
\end{tabular}
\end{table}

\subsection*{Apéndice A. Resultados del Algoritmo Genético por corrida}
\addcontentsline{toc}{subsection}{Apéndice A. Resultados del Algoritmo Genético por corrida}

\begin{table}[H]
\centering
\caption{Configuración base (población 100, mutación $1/n$): resultado de cada corrida}
\label{tab:ap_base}
\begin{tabular}{rrrrr}
\toprule
Semilla & Costo & Antenas & Generaciones & Tiempo (s) \\
\midrule
0 & 48\,878 & 22 & 104 & 2.5 \\
1 & 52\,945 & 25 & 96 & 1.8 \\
2 & 53\,807 & 24 & 93 & 1.9 \\
3 & 51\,793 & 24 & 96 & 2.1 \\
4 & 52\,007 & 24 & 98 & 2.1 \\
5 & 50\,253 & 23 & 103 & 2.1 \\
6 & 49\,501 & 22 & 96 & 2.0 \\
7 & 52\,551 & 24 & 93 & 2.0 \\
8 & 51\,881 & 24 & 94 & 2.1 \\
9 & 52\,603 & 24 & 93 & 2.2 \\
\midrule
Mejor / media / peor & \multicolumn{4}{r}{48\,878 / 51\,622 / 53\,807} \\
\bottomrule
\end{tabular}
\end{table}
\begin{table}[H]
\centering
\caption{Configuración extendida (población 200, mutación $2/n$): resultado de cada corrida}
\label{tab:ap_ext}
\begin{tabular}{rrrrr}
\toprule
Semilla & Costo & Antenas & Generaciones & Tiempo (s) \\
\midrule
100 & 50\,209 & 22 & 449 & 16.6 \\
101 & 50\,253 & 23 & 617 & 23.5 \\
102 & 50\,403 & 23 & 423 & 15.8 \\
103 & 50\,681 & 23 & 421 & 15.7 \\
104 & 50\,669 & 23 & 425 & 16.0 \\
105 & 50\,253 & 23 & 416 & 15.2 \\
106 & 49\,678 & 23 & 418 & 15.6 \\
107 & 50\,253 & 23 & 419 & 16.0 \\
108 & 49\,920 & 23 & 421 & 15.6 \\
109 & 52\,038 & 24 & 415 & 15.7 \\
\midrule
Mejor / media / peor & \multicolumn{4}{r}{49\,678 / 50\,436 / 52\,038} \\
\bottomrule
\end{tabular}
\end{table}

\subsection*{Apéndice B. Mejor solución entera del método exacto}
\addcontentsline{toc}{subsection}{Apéndice B. Mejor solución entera del método exacto}
HiGHS terminó por límite de tiempo (600 s), por lo que esta solución es la mejor encontrada y no está certificada como óptima: costo 50\,119, cota inferior 30\,817 y brecha 38.51\,\%. Las 22 antenas seleccionadas (numeradas desde 1) son:

\begin{quote}
13, 19, 44, 45, 73, 79, 94, 200, 210, 277, 307, 334, 357, 366, 401, 404, 411, 432, 444, 450, 473, 475.
\end{quote}

\subsection*{Apéndice C. Código fuente}
\addcontentsline{toc}{subsection}{Apéndice C. Código fuente}
Se incluyen los tres módulos centrales. El código completo (experimentos, figuras e interfaz) está en el repositorio indicado en el Apéndice D.

\lstinputlisting[caption={Carga de datos (\texttt{setcover\_datos.py})}]{codigo/setcover_datos.py}
\lstinputlisting[caption={Solución exacta con HiGHS (\texttt{solucion\_exacta.py})}]{codigo/solucion_exacta.py}
\lstinputlisting[caption={Algoritmo Genético con reparación (\texttt{algoritmo\_genetico.py})}]{codigo/algoritmo_genetico.py}

\subsection*{Apéndice D. Reproducibilidad y repositorio}
\addcontentsline{toc}{subsection}{Apéndice D. Reproducibilidad y repositorio}
Todo el código, los datos y los resultados están publicados en GitHub: \url{https://github.com/IsabelVargas2/set-cover-antenas}. Para reproducir los resultados:
\begin{enumerate}
\item Instalar las dependencias: \texttt{pip install -r requirements.txt}.
\item Resolver con el método exacto: \texttt{python correr\_exacto.py}.
\item Ejecutar los experimentos del Algoritmo Genético: \texttt{python experimentos\_ga.py} y \texttt{python experimento\_extendido.py}.
\item Generar las figuras: \texttt{python hacer\_figuras.py}.
\item Abrir la interfaz visual: \texttt{streamlit run app.py}.
\end{enumerate}
Los tiempos dependen del equipo; los costos del Algoritmo Genético dependen de la semilla, que se fija en cada corrida para que los resultados se puedan repetir.
